\title{Large Language Models Can Automatically Engineer Features for Few-Shot Tabular Learning}

\begin{abstract}
Large Language Models (LLMs) possess impressive skills in addressing complex and novel reasoning tasks, suggesting a strong aptitude for tabular learning, which underpins many practical scenarios. This paper introduces \textsf{FeatLLM}, a new in-context learning methodology that utilizes LLMs as feature generators, restructuring the input dataset to maximize tabular predictive accuracy. These constructed features are subsequently leveraged by a straightforward downstream model, such as linear regression, to estimate class probabilities, enabling robust few-shot performance. Notably, \textsf{FeatLLM} restricts the inference phase to this predictive model operating solely on the engineered features. Distinct from current LLM-powered solutions, our approach dispenses with the necessity to consult the LLM for each individual inference-time input. The method requires nothing more than API-level interaction with LLMs and is not impeded by prompt length restrictions. Our extensive experiments on numerous tabular datasets spanning multiple fields demonstrate that \textsf{FeatLLM} develops superior rule sets, achieving average performance gains of 10\% over methods such as TabLLM and STUNT.
\end{abstract}

\section{Introduction}
Large Language Models (LLMs) have recently exhibited substantial capacity to handle new tasks by generating appropriate outputs without dedicated task-specific adjustment~\cite{creswell2022selection,imani2023mathprompter,taylor2022galactica}. With their pretraining on massive textual datasets, these models internalize significant background knowledge, reducing reliance on subject-matter experts across a variety of domains~\cite{Genomes2021,singhal2023large,wilcox2003role}. This attribute is pivotal in automating numerous functions, including tasks like dialogue interpretation~\cite{finch2023leveraging} and automatic program correction~\cite{fan2023automated}. Moreover, LLMs' understanding is enhanced when task instructions and several illustrative examples are provided within a prompt, allowing them to discern contextual cues and focus on salient variables and instances~\cite{brown2020language,zhao2023survey}. These attributes underscore LLMs' proficiency in data analysis and signify their utility as advanced feature constructors.

The capacity of LLMs for encoding background information and logical reasoning is increasingly being harnessed within tabular data analysis. For example, contextual knowledge such as the predisposition of older age groups to specific diseases is valuable for clinical predictive models. Approaches in this space commonly express the task and feature characteristics in natural language, transform tabular inputs for compatibility, and subsequently process this serialized data through the LLM~\cite{dinh2022lift,wang2023anypredict}. These methods typically adopt either in-context learning paradigms~\cite{nam2023semi} or lightweight fine-tuning strategies~\cite{dinh2022lift,hegselmann2023tabllm}. Empirical evidence shows that tapping into the prior knowledge encoded in LLMs enables these systems to consistently surpass traditional tabular learning methods, particularly when only limited training data is available~\cite{hegselmann2023tabllm}.

Existing methods for tabular learning utilizing large language models (LLMs) encounter several drawbacks. One major issue is that end-to-end prediction requires at least one LLM inference per data instance, resulting in high computational overhead. Additionally, achieving high predictive accuracy frequently depends on fine-tuning the LLM, which becomes impractical for LLMs that lack openly available training infrastructure or tuning APIs. This is especially problematic for recently introduced LLMs that are only accessible through restricted inference APIs~\cite{anil2023palm,openai2023gpt4}. Further, most current techniques do not scale effectively with long prompts~\cite{liu2023lost}. When tabular datasets contain many features, resulting in prompts that exceed manageable lengths, methods often become unusable or their performance suffers. Collectively, these limitations hinder the practical use of LLM-driven tabular learning in real-world contexts.

The methodology we introduce reframes the conventional end-to-end predictive use of LLMs. Rather than using LLMs solely to predict outcomes, we exploit their feature engineering capabilities, allowing them to offer a novel pathway for deploying their domain expertise. We propose a new in-context learning paradigm, \textsf{FeatLLM}, designed to elucidate the ``criteria'' that inform LLM decision-making. For example, in disease prediction scenarios, an LLM can deduce and articulate rules mapping specific feature patterns to disease identification. These extracted criteria serve as the basis for constructing new features that replace those originally present, which, in turn, streamlines subsequent modeling steps. This paradigm assigns the LLM the responsibility of feature creation, permitting the use of simpler downstream models—once the informative features are extracted, complex inference is unnecessary, thereby reducing latency. In-context learning capabilities are leveraged to extract these features without training the LLM, relying on prompt demonstrations instead. Moreover, ensemble-inspired methods such as feature bagging~\cite{breiman2001random} can be employed to address issues related to excessive prompt length, thus mitigating one of the chief obstacles to scalability.

\textsf{FeatLLM} decomposes the prompt for feature engineering into two distinct subtasks: (1) comprehending the problem at hand and deducing the associations between features and targets; and (2) using insights from the initial step and provided few-shot instances to formulate precise decision rules for class discrimination. This form of structured reasoning encourages LLMs to disregard irrelevant features when constructing rules. The resulting rules are subsequently implemented on the data (interpreted as program code), converting the dataset into binary indicators reflecting sample adherence to each rule. Ultimately, a simple model is trained to infer class probabilities using these binary features. Model reliability is further enhanced through an ensemble strategy, which involves recurrent feature generation paired with feature bagging. By tuning the counts of features and samples involved in bagging, the issue of overly large prompts is alleviated. \textsf{FeatLLM} avoids training requirements and incurs minimal inference overhead, facilitating practical adoption of LLMs for diverse real-world applications.

Our assessment of \textsf{FeatLLM} spans 13 tabular datasets under low-shot conditions, revealing strong and consistent performance. We demonstrate that LLMs effectively leverage prior domain knowledge together with dataset-derived insights to produce informative features. The proposed framework consistently surpasses leading few-shot learning baselines under various scenarios. The accompanying code is publicly available through an anonymized GitHub repository at~\texttt{https://github.com/Sungwon-Han/FeatLLM}.

\section{Related Work}

\subsection{Few-Shot Learning with Tabular Data} Recent progress in few-shot learning has allowed models to derive robust and transferable features from data, even when labeling resources are scarce~\cite{chen2018closer}. Originally centered on image-based problems, few-shot methodologies have been demonstrated to exhibit considerable versatility, now encompassing modalities such as text, audio, and, more recently, tabular datasets~\cite{majumder2022few,nam2023stunt,schick2021exploiting}. Nevertheless, the challenge with extremely limited labeled examples is that models can latch onto non-generalizable patterns~\cite{bartlett2021deep}. As a result, there has been a shift toward leveraging extra sources of information or embedding task-specific priors to steer learning towards more reliable representations. Strategies include exploiting abundantly available unlabeled data—combined with targeted augmentation techniques—for semi-supervised paradigms~\cite{ucar2021subtab,yoon2020vime}, as well as utilizing broad, unsupervised pre-training for improved model initialization~\cite{somepalli2021saint}. These unsupervised schemes might require occluding or corrupting sections of the data, subsequently using reconstruction as a learning signal~\cite{arik2021tabnet,majmundar2022met}, or devising contrastive tasks via tailored augmentation~\cite{bahri2022scarf,chen2020simple}. However, the intrinsic heterogeneity found in tabular datasets complicates the discovery of augmentation methods that transfer universally, and empirical evidence has shown that these methods yield limited gains in low-data regimes~\cite{nam2023stunt}.

To counter these difficulties, more recent frameworks propose first exposing models to a wide collection of heterogeneous tabular data that are not necessarily associated with the downstream target task, followed by knowledge transfer to the task at hand~\cite{zhang2023generative,levin2022transfer,zhu2023xtab}. A representative example is TransTab~\cite{wang2022transtab}, which adopts transformers that exploit semantic clues from column names, in addition to raw cell entries, for developing general-purpose column relationships. Such acquired cross-table representations underpin successful zero-shot and few-shot inference for newly encountered table tasks. In another instance, TabPFN~\cite{hollmann2022tabpfn} pre-trains a Bayesian neural network on synthetic datasets simulated from diverse distributions. Upon pre-training, predictions on new tasks can be made directly using both support and query data, without need for additional model updating. The transfer of inductive biases learned from these varied datasets leads to performance superior to classical decision-tree approaches and has demonstrably enhanced few-shot task accuracy.

\subsection{Language-Interfaced Tabular Learning} The integration of large language models (LLMs) presents a novel pathway to inject prior knowledge for tabular learning~\cite{anil2023palm,openai2023gpt4}. Due to their extensive exposure to vast and varied textual resources, LLMs are capable of performing well on tasks outside their training distribution and have broad applicability in different fields~\cite{brown2020language}. Contemporary research explores channeling the information and knowledge captured by LLMs into tabular applications. The standard approach entails converting tabular datasets into textual sequences, which are then used as part of the LLM's prompt~\cite{dinh2022lift,hegselmann2023tabllm,wang2023anypredict}. By enriching the input with both raw table content, supplementary feature documentation, and explicit task descriptions, these systems can infer the dependencies and associations necessary for task-specific predictions. Methodological advances in this space range from tuning model parameters efficiently through techniques like LoRA~\cite{hu2021lora} and IA3~\cite{liu2022few}, to applying in-context learning via prompt-based demonstration with only a few labeled examples—requiring no gradient-based training—~\cite{dinh2022lift,hegselmann2023tabllm,nam2023semi,slack2023tablet}.

While previous approaches have proven useful for few-shot tabular learning, their reliance on routing inference through the LLM for every sample introduces practical limitations for deployment in industrial settings. The current work seeks to shift away from the standard black-box paradigm, wherein predictions are handled end-to-end by an LLM. Instead, it proposes utilizing the LLM to explicitly generate the underlying rules or criteria that define each class. \textsf{FeatLLM} interprets these rules by converting them into programmatic features, thereby facilitating prediction without repeated LLM inference and using in-context learning to distill rules based on both prior knowledge and the provided few-shot examples.

\section{Methods}
The objective of \textsf{FeatLLM} is to identify ``rules''—that is, the specific conditions distinguishing each possible class—from the handful of provided samples, as visualized in Figure~\ref{fig:main_model}. These class-specific rules are extracted with an LLM, parsed, and transformed into binary feature indicators for any input sample, reflecting whether each rule is met. These binary features then serve as input to a dot product operation alongside non-negative, learnable weights to compute each class's probability. Model training iteratively adjusts the weights to reflect the contribution of individual features to the classification outcome (Section~\ref{Sec:training}). The methodology employs repeated sampling via bagging and aggregates predictions through ensembling. The subsequent sections elaborate on the problem setup and describe each stage in detail.

\vspace*{-0.12in}
\paragraph{Problem formulation.} We begin with a tabular dataset containing $N$ labeled examples, defined as $\mathcal{D} = \{(\mathbf{x}^i, \mathbf{y}^i)\}_{i=1}^{N}$. Each instance $\mathbf{x}^i$ is represented as a $d$-dimensional feature vector, and the dataset is accompanied by natural language identifiers for each feature, $F = \{f_j\}_{j=1}^{d}$, with individual feature descriptions provided separately. Within a classification context, the target $\mathbf{y}^i$ refers to a member of a finite set of predefined categories. For $k$-shot learning scenarios, the model is trained on a randomly selected subset of $k$ ($k < N$) labeled examples.

\subsection{Prompt Design for Extracting Rules}
\label{Sec:prompt_design}
In order to assist the LLM in deriving rules through a structured and informed reasoning procedure, we prompt the model to emulate the approach taken by a skilled human expert when performing tabular analysis. The constructed prompt encompasses the following three essential elements (illustrated in Fig.~\ref{fig:prompt_for_rule} and detailed in Appendix~\ref{sec:example_rule}):

\vspace*{-0.12in}
\paragraph{Basic information description.} This segment supplies the foundational content needed for task resolution (highlighted as orange in Fig.~\ref{fig:prompt_for_rule}). It presents a task statement with associated label details, feature descriptions, and sample demonstrations. The task is articulated as a question (such as ``Does this patient's myocardial infarction complications data indicate chronic heart failure? Yes or no?"), with further examples accessible in Appendix~\ref{sec:task_desc}. The feature descriptions specify the data type and, where relevant, provide the feature definition. For features that are categorical, sample category values may also be listed. Example demonstrations comprise a selection of training examples reformulated as textual data alongside their verified labels. The serialization process adopts conventions established by previous works~\cite{dinh2022lift,hegselmann2023tabllm}:

\begin{align}
&\text{Serialize}(\mathbf{x}^i,\mathbf{y}^i, F) = \nonumber \\ 
&``\ f_1\ \text{is}\ \mathbf{x}^i_1.\ ... \ f_d\  \text{is}\  \mathbf{x}^i_d.\ \ \text{Answer:}\  \mathbf{y}^i\ ”
\end{align}

\vspace*{-0.12in}
\paragraph{Reasoning instruction.} Our objective is to improve the reasoning capabilities of the LLM by supplying explicit reasoning directives (see the blue highlights in Fig.~\ref{fig:prompt_for_rule}). The reasoning guidance starts with a sentence modeled on the chain-of-thought method~\cite{wei2022chain}. Subsequently, the inference procedure is partitioned into two distinct phases. In the first phase, the LLM is prompted to hypothesize causal influences or general tendencies linking features and task descriptions. This reasoning takes place independently from example demonstrations, instead utilizing broad general knowledge and common sense, empowering the LLM to draw on its prior understanding of the topic. In the second phase, the LLM incorporates sample demonstrations together with the previously established knowledge to infer classification rules. This dual-step process mitigates the risk of the model being misled by irrelevant feature associations and encourages concentration on pertinent attributes. To prevent both overly sparse rule sets and unnecessarily lengthy prompts—which could result in underfitting or excessive complexity—a constant number of rules (namely, 10) is specified for extraction per class.\footnote{Section~\ref{sec:analysis} provides an evaluation regarding rule quantity.} Moreover, the LLM is advised to consult the feature type information present in the preliminary descriptions while deriving rules, thereby minimizing the chance of generating rules with mismatched types.

\vspace*{-0.12in}
\paragraph{Response instruction.} To promote clarity in the parsing and subsequent application of the derived rules, we furnish the LLM with detailed response formatting instructions (illustrated in yellow within Fig.~\ref{fig:prompt_for_rule}). The prompt outlines the desired format for answers in each class (i.e., response formatting) as well as the requisite template for individual rules (i.e., rule formatting). Such directives enable the LLM to adapt rule structure in accordance with feature type. When no additional instructions are provided, the LLM may combine several rules using logical connectors such as AND or OR. An illustrative outcome is presented in Appendix~\ref{sec:outcome-rule}.

\subsection{Inferring Class Likelihood via Rules}
\label{Sec:training}

\paragraph{Parsing rules for feature generation.} The rules obtained in the prior step are employed to construct novel binary features for each class. These features capture whether a particular sample adheres to the class-specific rules (yielding a value of either '0' or '1'). Such features serve as indicators when estimating the likelihood of class membership for a given sample. Because these rules originate from the LLM and are described in natural language, an automatic parsing methodology is necessary to facilitate feature creation. Although formatting constraints are imposed during rule generation to streamline parsing, the LLM's responses may still exhibit irregular syntax~\cite{kovavc2023large}. For example, expressions like ``$>$” or ``$<$” might be replaced by verbal forms such as ``greater than" or ``less than," increasing parsing complexity.

To overcome difficulties posed by noisy text, rather than implementing elaborate parsing logic, we harness the LLM itself. The LLM possesses robust semantic parsing capabilities and can synthesize simple code from textual instructions (due to its code-centric training corpus). To capitalize on this, we construct prompts containing the target function name, descriptions of inputs and outputs, as well as the inferred rules, and then present these to the LLM. Figures~\ref{fig:prompt_for_parsing} and ~\ref{sec:example_parsing}-\ref{sec:example_parse_outcome} in the Appendix illustrate example prompts and parsing outcomes. The Python \texttt{exec()} function is utilized to execute the code produced by the LLM, facilitating the transformation of data.

\vspace*{-0.12in}
\paragraph{Inferring class likelihood.} With new binary rule-based features for each class, one rudimentary technique for estimating a sample's likelihood of belonging to a class is to count the number of satisfied rules—equivalently, sum the binary features corresponding to each class. Nonetheless, the influence of individual rules may differ, thus necessitating the determination of their relative importance from training data. To achieve this, we apply a standard machine learning (ML) model that can assign learned weights to the binary features for each class. For instance, a linear model with zero bias may be used. Let the binary features for sample $\mathbf{x}^i$ and class $k$ be denoted by $\mathbf{z}_k^i \in \{0, 1\}^R$, where $R$ signifies the number of rules for a class and $c$ denotes the total number of classes. The class probability vector $\mathbf{p}^i$ for sample $\mathbf{x}^i$ is then determined as follows:

\begin{align}
\text{logit}_k^i &= \max(\mathbf{w}_k, 0) \cdot \mathbf{z}_k^i, \nonumber \\
\mathbf{p}^i &= \text{Softmax}({[\text{logit}_{1}^i, ..., \text{logit}_{c}^i]}). \label{eq:infer_prob}
\end{align}

In this context, $\mathbf{w}_k$ denotes the parameters of the linear model that are trained, with each weight indicating the relevance of the corresponding feature. By restricting these weights to a positive domain, we ensure that the LLM-generated features cannot decrease the class probability during the learning phase. The output logits are mapped to class probabilities via the softmax operation. Training involves minimizing the cross-entropy loss to learn the weights $\mathbf{w}_k$, using a small number of labeled training samples.

\vspace*{-0.12in}
\paragraph{Ensembling with bagging.} To generate the final prediction, we repeat this procedure multiple times, resulting in several inference models whose outputs are aggregated. With the set of learned weights from $T$ independent runs, $\{\mathbf{w}[t]\}_{t=1}^T$, class probabilities $\mathbf{p}^i$ are determined for each model and then averaged as shown in Eq.~\ref{eq:infer_prob}.

To foster diversity in rule generation for the ensemble, we employ an increased temperature during LLM inference. We further shuffle the sequence of few-shot exemplars based on evidence that the order alters the output of the LLM~\cite{lu2022fantastically}\footnote{Further discussion on rule diversity appears in Appendix~\ref{sec:diversity}}. To take advantage of diverse predictions and strengthen robustness, bagging is implemented by randomly selecting subsets of features or samples for each repetition, especially beneficial when dealing with a larger pool of features or training data. This ensemble strategy confers two primary benefits: firstly, any erroneous rules produced by the LLM in individual repetitions are balanced by accurate rules from others, thereby boosting the framework's reliability. This aligns with LLM self-consistency~\cite{wang2022self}, since rules seen across multiple runs tend to be correct. Secondly, the ensemble technique mitigates the issue of limited prompt size when dealing with extensive tabular datasets—bagging allows prompt reduction and maintains performance. While a single rule set might overlook certain features, combining many weak learners from numerous trials helps compensate for gaps in coverage and counteracts the limitations faced by individual models.

\section{Experiments}
We assess the effectiveness of \textsf{FeatLLM} using a broad range of tabular datasets, accompanied by a series of in-depth analyses to elucidate the operational aspects of our framework. To accommodate space limitations, supplementary experiments—including tests with alternative LLMs, qualitative discussions, and other evaluations—are provided in the Appendix.

\subsection{Performance Evaluation}
\paragraph{Datasets.}
Our study employs 13 distinct datasets covering binary and multi-class classification tasks: (1) Adult~\cite{asuncion2007uci} focuses on identifying individuals earning above \$50,000 per year; (2) Bank~\cite{moro2014data} predicts customer subscription to term deposits; (3) Blood~\cite{yeh2009knowledge} anticipates repeat blood donation events; (4) Car~\cite{kadra2021well} evaluates automobile quality; (5) Communities~\cite{misc_communities_and_crime_183} estimates regional crime rates; (6) Credit-g~\cite{kadra2021well} determines whether a person constitutes a favorable or unfavorable credit risk; (7) Diabetes\footnote{\texttt{kaggle.com/datasets/uciml/pima-indians-diabetes-database}} diagnoses diabetes for individuals; (8) Heart\footnote{\texttt{kaggle.com/datasets/fedesoriano/heart-failure-prediction}} predicts the likelihood of coronary artery disease occurrence; (9) Myocardial~\cite{misc_myocardial_infarction_complications_579} detects chronic heart failure presence.

Further, we include both contemporary and synthetic datasets, which are less commonly analyzed in prior work and were absent from LLM pre-training corpora: (10) Cultivars forecasts soybean cultivar grain yield\footnote{\texttt{archive.ics.uci.edu/dataset/913}}; (11) NHANES categorizes individuals into age brackets based on their records\footnote{\texttt{archive.ics.uci.edu/dataset/887}}; (12) Sequence-type classifies sequences as arithmetic, geometric, Fibonacci, or Collatz; (13) Solution-mix determines whether a mixture composed of four solutions surpasses a concentration threshold of 0.5. The Cultivars and NHANES datasets were recently contributed to the UCI Machine Learning Repository post-September 2023, confirming their exclusion from the pre-training data of the LLM API (gpt-3.5-turbo-0613) leveraged by our framework. Sequence-type and Solution-mix are synthetic, negating any possibility of memorization by the LLM API during evaluation.

Dataset properties such as sample size and feature complexity are outlined in Table~\ref{Tab:basic-desc}. Each dataset includes explicit nomenclature and descriptions for all attributes. High-dimensional datasets like `Communities' and `Myocardial,' which contain more than 100 attributes, present additional difficulties linked to LLM context window limitations.

\vspace*{-0.12in}
\paragraph{Baselines.}
Comparative analyses were performed against nine baseline methods. Three conventional models serve as initial baselines: (1) Logistic Regression (LogReg), (2) XGBoost~\cite{chen2016xgboost}, and (3) RandomForest~\cite{ho1995random}. Two semi-supervised approaches, suitable for settings with unlabeled data, include (4) SCARF~\cite{bahri2022scarf} and (5) STUNT~\cite{nam2023stunt}, though it is important to note that substantial amounts of unlabeled data may not be accessible in practical scenarios. The following baseline, (6) TabPFN~\cite{hollmann2022tabpfn}, is characterized by synthetic pre-training with diverse data distributions. The remaining approaches harness LLMs: (7) In-context learning (In-context)~\cite{wei2022emergent}, (8) TABLET~\cite{slack2023tablet}, and (9) TabLLM~\cite{hegselmann2023tabllm}. In-context learning places few-shot examples directly in the input prompt, omitting further model adaptation. TABLET augments prompts with supplementary guidance obtained from external classifiers via rule-based sets and prototype information, aiming to refine inference. TabLLM employs parameter-efficient tuning—specifically IA3~\cite{liu2022few}—allowing LLM adaptation through a limited number of training examples.

\vspace*{-0.12in}
\paragraph{Implementation details.}
\textsf{FeatLLM} is built on GPT-3.5 as its primary LLM, although its architecture accommodates any LLM, as demonstrated with various backbones in Appendix~\ref{sec:backbones}. The LLM is queried with a temperature of 0.5, and the API’s top-p parameter remains at 1, its standard value. For the ensemble setup, 20 models are used, and feature extraction relies on generating 10 distinct rules. The effects of different hyper-parameter values are presented in Figure~\ref{fig:hyperparameter2} and discussed further in Appendix~\ref{sec:hyper-parameter}. Training the linear model employs the Adam optimizer with a learning rate of 0.01 over 200 epochs, leveraging $k$-fold cross-validation to pinpoint the best epoch. When reproducing baselines, in-context learning utilizes GPT-3.5, while TabLLM leverages T0~\cite{sanh2022multitask} in accordance with its original setup; this is due to TabLLM's limitation to only LLMs with open checkpoints, excluding GPT-3.5. For conventional algorithms like LogReg, XGBoost, and RandomForest, parameter selection is accomplished via Grid-search, alongside $k$-fold cross-validation, with $k$ set at 2 or 4 so that each training partition contains samples from every class. Additional implementation specifics are documented in Appendix~\ref{sec:baselines}.

\vspace*{-0.12in}
\paragraph{Results.}
Tables~\ref{tab:main1} and \ref{tab:main2} show comparative performance evaluations across several datasets. Each experiment is conducted three times, using different splits of training and testing sets, and we report both the mean and the standard deviation of the AUC scores. \textsf{FeatLLM} emerges as either the top-performing or second-best method among all baseline models examined. Remarkably, our approach attains results close to those of traditional tabular learning baselines using 32 shots, while requiring only 4 shots (as depicted in Fig.~\ref{fig:compares}). Although models such as STUNT and TabPFN exhibited noticeable gains on select datasets, they did not yield consistently stronger improvements relative to standard machine learning algorithms. Additionally, LLM-based techniques—including in-context learning, TABLET, and TabLLM—struggled to conduct inference on tabular data sets with a large feature space. Our proposed framework, which leverages bagging and ensemble strategies, maintains robust effectiveness even as the number of features grows, achieving strong results. It is worth mentioning that the ensemble and bagging methodology employed here can be integrated into pre-existing LLM-based solutions. However, this modification would double inference requirements per sample, thereby considerably increasing computational demands and rendering the process less practical.

\subsection{Analysis \& Discussion}
\label{sec:analysis}
\paragraph{Ablation study.} We examine the impact of each primary framework component through a set of ablation experiments, specifically analyzing: (1) \textsf{-Tuning}: excluding the process of optimizing linear model weights, instead using the sum of binary features to produce the logit; (2) \textsf{-Ensemble}: removing ensemble aggregation; (3) \textsf{-Description}: eliminating feature descriptions; and (4) \textsf{-Reasoning}: bypassing the initial step in the reasoning instructions used for rule generation.

Table~\ref{tab:ablation} reports the mean change in AUC following these ablations, aggregated across all datasets. Each component’s removal consistently leads to lower performance. Of particular note, omitting weight tuning has a growing detrimental effect as the number of shots increases, suggesting that access to more data enables more accurate learning of rule importance, thereby heightening the value of the tuning step. Conversely, the absence of feature descriptions and reasoning instructions incurs greater performance drops when the number of shots is limited. This implies that leveraging the prior knowledge within the LLM is especially valuable under data scarcity. Finally, applying the ensemble method persistently results in performance improvements under all test conditions.

\vspace*{-0.12in}
\paragraph{Training \& inference time comparison.} We evaluate the training and inference times across various techniques. Our experiments are conducted on the Adult dataset, utilizing a training set with 16 shots. For most methods, we use a single A100 GPU by default, though TabLLM relies on four A100 GPUs to enable model parallelism. Table~\ref{tab:cost} presents detailed runtime metrics. Significantly, \textsf{FeatLLM} achieves relatively low inference latency, on par with traditional approaches. Conversely, inference times for other LLM-based methods—including In-context learning, TABLET, and TabLLM—are considerably greater. With respect to training time, \textsf{FeatLLM} aligns closely with other LLM-driven approaches; it requires just 30 API queries, which are feasible within most budgets and amenable to parallel execution.

\vspace*{-0.12in}
\paragraph{Quality of features generated by \textsf{FeatLLM}.}
To analyze the effectiveness of rule-based features produced by \textsf{FeatLLM} for real-world tasks, we run a simple assessment. We compute the mean AUROC between the newly constructed features and the target labels, subsequently determining the proportion of features whose AUROC exceeds 0.5 for each dataset. The outcome is summarized in Table~\ref{tab:quality}. The findings indicate that most generated rules are substantially connected to the target variable, offering valuable contributions for task performance (see column ``Before tuning"). Additionally, the linear model tuning procedure serves to filter out less informative rules—those with learned weights below a certain threshold—automatically (see column ``After tuning"). The threshold, set at 0.1, is selected based on the global weight distribution histogram.

\vspace*{-0.12in}
\paragraph{Effect of adding spurious correlations.} To evaluate the ability of different approaches to eliminate spurious correlations—especially under limited labeled data—we performed an experiment. In this experiment, the Heart dataset served as the main task, and we appended random columns from the Adult dataset that bear no relevance to the Heart dataset. We systematically varied the number of these extraneous columns and tracked changes in model accuracy as more irrelevant information was introduced. To ensure a fair assessment, all methods were supplied with exactly 8 labeled examples, omitting access to any unlabeled data. This design choice excludes methods such as SCARF and STUNT from the comparison.

Figure~\ref{fig:spurious} illustrates how model performance evolves as spurious correlations are incrementally added. In comparison to standard baselines, \textsf{FeatLLM} experienced the smallest drop in accuracy as irrelevant features increased. The resilience of our framework can be attributed to its mechanism: when deriving rules, it encourages alignment between feature relevance and prior knowledge about the task, which restricts the inclusion of rules dependent on noise variables, thus promoting robustness. Notably, less than 2\% of the final rules in \textsf{FeatLLM} involved the appended unrelated columns. On the other hand, \textsf{FeatLLM} can still erroneously capture misleading associations, particularly if randomly appended columns happen to originate from datasets with inherent similarity to the Heart dataset (see Appendix~\ref{sec:spurious-appendix}).

\vspace*{-0.12in}
\paragraph{Hyper-parameter analysis}
\textsf{FeatLLM} is governed by two primary hyper-parameters: the ensemble count and the number of rules per class retrieved in each LLM query. To assess the influence of these hyper-parameters on model performance, we ran systematic experiments. Our findings reveal that increasing the ensemble count leads to enhanced performance; however, improvement plateaus when the ensemble size approaches 20 (refer to Fig.~\ref{fig:appendix-hyperparameter1} in the Appendix). Regarding the number of rules extracted, no statistically significant differences emerged, but we identified 10 as the optimal number, striking a balance between prompt length and prediction efficacy (see Fig.~\ref{fig:hyperparameter2}). This outcome may be attributed to the expansion of input dimensionality with more rules, providing additional information but also increasing the risk of overfitting, particularly in low-shot settings.

\section{Conclusion}
We introduce \textsf{FeatLLM}, which elevates LLM-driven approaches for few-shot tabular learning. \textsf{FeatLLM} distinguishes itself by leveraging LLMs not merely for direct sample prediction, but for synthesizing predictive criteria as a form of feature engineering. By integrating rule generation with a bagging-based ensemble methodology, \textsf{FeatLLM} ensures low inference overhead, functions solely with API-level access (requiring no retraining), and circumvents limitations related to tabular feature dimensionality. The proposed method demonstrates robust predictive results and presents a compelling solution for real-world tabular applications that demand data efficiency.

\textsf{FeatLLM} is specifically tailored for scenarios with limited sample availability. Looking ahead, our intention is to broaden its scope to facilitate feature generation for datasets containing larger sample sizes, to investigate more varied feature constructions beyond rule extraction, and to promote interpretability, thus making the technique applicable to a wider range of tasks.

\section{Broader Impact}
The outlined \textsf{FeatLLM} framework leverages the prior knowledge encapsulated within LLMs for enhanced tabular learning performance, surpassing conventional supervised learning approaches. By advancing data efficiency and supplementing limited training data with the pretrained insights of LLMs, our framework offers a means to diminish overfitting. \textsf{FeatLLM} holds significant promise for professionals in domains such as finance or healthcare, where acquiring labeled data is often costly or infeasible, and pretrained domain-relevant knowledge from LLMs can be advantageous. A critical future step for widespread adoption includes evaluating the implications of inherent societal biases and misinformation present in LLMs’ prior knowledge, which may substantially shape prediction outputs and even override evidence from labeled data in practical deployments.

\end{document}